\documentclass[11pt,letterpaper]{article}
\usepackage[T1]{fontenc}
\usepackage{lmodern,microtype}
\usepackage[margin=1in]{geometry}
\usepackage{amsmath,amssymb,amsthm}
\usepackage{booktabs,tabularx}
\usepackage{xcolor}
\usepackage[colorlinks=true,linkcolor=blue!40!black,citecolor=blue!40!black,
            urlcolor=blue!40!black]{hyperref}

\newcommand{\manuscripttitle}{Detecting signed powers over fields,\\with an application to hidden binary wirings}
\newcommand{\manuscriptkeywords}{detecting matrix, signed power, algebraic number, inversion count, coordinate permutation}
\hypersetup{
  pdftitle={Detecting signed powers over fields, with an application to hidden binary wirings},
  pdfauthor={Joseph M. Shunia},
  pdfsubject={Detecting matrices for signed powers and binary coordinate permutations},
  pdfkeywords={\manuscriptkeywords}
}
\newtheorem{theorem}{Theorem}[section]
\newtheorem{lemma}[theorem]{Lemma}
\newtheorem{proposition}[theorem]{Proposition}
\newtheorem{corollary}[theorem]{Corollary}
\numberwithin{equation}{section}
\setlength{\emergencystretch}{1.5em}
\addtolength{\skip\footins}{6pt}

\title{\vspace{-1.5\baselineskip}%
  \manuscripttitle}
\author{Joseph M. Shunia\thanks{Independent researcher, Ann Arbor, Michigan, United States.
  Email: \href{mailto:jshunia@gmail.com}{\texttt{jshunia@gmail.com}}.}}
\date{November 22, 2026}

\begin{document}
\maketitle
\vspace{-\baselineskip}

\begin{abstract}
We determine the number of signed-sum queries required to recover $n$ unknown entries that are
zero or signed powers of a known field element. For any fixed substitution fraction below one
quarter, a transcendental radix or a root of unity in characteristic zero requires order
$n/\log n$ queries. An algebraic radix of infinite order requires order $n\log\log n/\log n$.
In odd characteristic the order is linear, and in characteristic two recovery is impossible
for $n\ge2$. The algebraic upper bound uses a height argument to shorten exponent gaps and
finite quotients that preserve large additive orders. As an application, we recover a hidden
permutation and complementation of binary coordinates from inversion counts at the same rate
as for powers of two.
\\[2mm]
{\bf 2020 Mathematics Subject Classification:} 94B25 (primary), 11R04, 05D40 (secondary).
\\[2mm]
{\bf Keywords:} \manuscriptkeywords.
\end{abstract}

\section{Introduction}\label{sec:intro}

The equality $\alpha^e+\alpha^e=\alpha^{e+1}$ holds for $\alpha=2$. If $\alpha$ is an
indeterminate, a sum of its powers instead records a separate coefficient at every exponent.
It is natural to ask how strongly this distinction affects recovery of a vector when each
observation is a single signed sum. The detecting matrices of Lindstr\"om \cite{lindstrom}
treat bounded integer alphabets. The noisy adder channel of Erdo\u{g}an, Maringer and
Polyanskii \cite{emp} already exhibits a one-quarter error threshold for binary signatures.
Here the exponent of even one unknown entry may be arbitrarily large.

Fix a field $K$ and $\alpha\in K^\times$. An unknown vector has $n$ independently chosen
coordinates in $D_\alpha=\{0\}\cup\{\pm\alpha^e:e\in\mathbb Z\}$. A query chooses one sign per
coordinate and returns their signed sum in $K$. A fixed $m$-row sign matrix $A$ produces the
word $Ax$. Answers are exact field elements, and cost is the number of queries. The radix and
field are known. We ask for unique recovery of every input after at most $\lfloor\eta
m\rfloor$ arbitrary answer substitutions. Write $M_\alpha(n,\eta)$ for the minimum positive
fixed length, and set it to infinity when recovery is impossible.

\begin{theorem}\label{thm:class}

For every fixed $0\le\eta<1/4$, the following table gives the order of $M_\alpha(n,\eta)$ as
$n$ grows. The field and radix are fixed in each row. Algebraicity and transcendence in
characteristic zero are over $\mathbb Q$.

\begin{center}\normalfont
\begin{tabularx}{\linewidth}{@{}Xl@{}}
\toprule
Field and radix & Number of queries \\
\midrule
Characteristic zero, transcendental radix or root of unity & $\Theta(n/\log n)$ \\
Characteristic zero, algebraic radix of infinite order & $\Theta(n\log\log n/\log n)$ \\
Odd characteristic, every radix & $\Theta(n)$ \\
Characteristic two, every radix & Impossible for $n\ge2$ \\
\bottomrule
\end{tabularx}
\end{center}

The lower orders hold even for deterministic adaptive noiseless queries. In every recoverable
case the supremum of asymptotically correctable fractions for fixed queries is $1/4$. The
table also holds with exponents restricted to nonnegative integers.
\end{theorem}

In the transcendental case every fixed matrix has the same minimum transcript distance as it
has on the ternary alphabet $\{-1,0,1\}$. The algebraic case requires a different reduction.
Lenstra's height estimate for lacunary polynomials \cite{lenstra} splits a vanishing sum into
shorter bands. We construct several finite quotients that together distinguish the shortened
powers and in which one repeated difference gives enough distinct signed sums.

For each fixed $\eta<1/4$ a single choice of infinite sign rows for each dimension has the
following property. Every fixed radix and field in a recoverable case eventually admits a
prefix of the indicated order. Prefix lengths and the first dimension for which they suffice
may depend on the radix or characteristic.

Section~\ref{sec:wiring} gives an integer-valued observation with exactly these measurement
rows. The number of inversions of a permuted and complemented binary word, after an output
mask is applied, is an affine function of a sum of signed powers of two. This provides both an
explicit decoder from $n$ noiseless counts and a smaller family of masks that corrects
substitutions.

\section{Characteristic and the finite-alphabet cases}\label{sec:elementary}

In characteristic two all signs equal one. Swapping the inputs $(1,0)$ and $(0,1)$ gives the
same answer to every query. In odd characteristic restrict first to the $2^n$ binary inputs.
Each query then has at most $p$ possible answers, even if chosen adaptively. This forces $m\ge
n\log 2/\log p$. For the upper bound draw independent fair signs. For any fixed nonzero
$v\in\mathbb F_p^n$, condition on every sign except one at a nonzero coordinate. Its two
outcomes differ, so a row vanishes on $v$ with probability at most $1/2$. By Hoeffding's
inequality \cite{hoeffding}, the probability of vanishing on at least $(1-2\eta)m$ of the rows
is at most $\exp[-2(1/2-2\eta)^2m]$. A union bound over fewer than $p^n$ vectors gives the
desired distance for $m=C_{p,\eta}n$ with exponentially small failure probability. For any
difference vector over $K$, take an $\mathbb F_p$-linear functional on the span of its
coordinates that is nonzero on one coordinate. Each nonzero projected answer certifies an
unequal original answer. This proves the upper bound even for unrestricted inputs in $K^n$.

In characteristic zero the binary inputs give $2^n$ possibilities and at most $2n+1$ answers
per query. For a fixed finite alphabet, the random row argument in Section~\ref{sec:count}
gives the matching $O(n/\log n)$ upper bound. This applies to roots of unity. For
transcendental $\alpha$, the Laurent monomials $\alpha^e$ are linearly independent over
$\mathbb Q$. Express each input as a list of ternary coefficient vectors, one for each
exponent. A row agrees on two inputs exactly when it agrees on every layer. Any two distinct
inputs differ on at least one layer. Their transcript distance is therefore at least the
minimum ternary distance. The ternary alphabet itself lies in $D_\alpha$, giving equality for
every matrix. Clearing a common monomial handles negative exponents.

For the error ceiling use the two attainable differences $(2,2,0,\ldots)$ and
$(2,-2,0,\ldots)$ of signed unit vectors with arbitrary identical remaining coordinates in
each pair. In characteristic different from two, exactly one gives a nonzero answer on each
row. The two distances sum to $m$, so one is at most $\lfloor m/2\rfloor$. Correcting $t$
substitutions requires distance at least $2t+1$. This proves the ceiling with the
finite-length rounding intact.

\section{The lower bound for an algebraic radix}\label{sec:lower}

Let $\alpha$ be algebraic of infinite order and put $E=\mathbb Q(\alpha)$, of degree $d$. The
companion matrix for multiplication by $\alpha$ in the power basis of $E$ has rational
entries. Write it $U/a$, where $a$ is a positive integer and $U$ is integral. For $0\le e<L$,
the power basis coordinates of $a^L\alpha^e$ are integers bounded by $C_\alpha^L$ in
magnitude. Indeed they are the entries of $a^{L-e}U^e$ applied to the first basis vector, and
one may take $C_\alpha\ge\max(2,a,\|U\|_\infty)$.

Take $L=\lfloor\log n/(\log\log n)^2\rfloor$ for sufficiently large $n$. The $L^n$ positive
power vectors are all distinct. After scaling by $a^L$, each answer belongs to an integer box
with at most $(2nC_\alpha^L+1)^d$ points. A depth-$m$ adaptive tree has at most that many
choices at each node, so
\[
m\ \ge\ \frac{n\log L}{d\log(2nC_\alpha^L+1)}
=\bigl(1/d-o(1)\bigr)\frac{n\log\log n}{\log n}.
\]

This part uses only bounded-degree arithmetic and gives the lower order for the noisy fixed problem as well.

\section{Shortening the gaps between exponents}\label{sec:gaps}

Use normalized absolute values $|\cdot|_v$ on $E$ extending the usual absolute values on
$\mathbb Q$, with local weights $w_v=[E_v:\mathbb Q_v]/d$. For $z\ne0$, the product formula
says $\prod_v |z|_v^{w_v}=1$. The archimedean weights sum to one. Put
\[
\lambda=\sum_v w_v\log^+|\alpha|_v .
\]

Here $\lambda>0$. If it were zero, $\alpha$ would be an algebraic integer with all conjugates
of absolute value at most one. The monic minimal polynomials of all its positive powers would
have uniformly bounded degrees and integer coefficients, by the elementary symmetric formulas.
There would be only finitely many possible powers, contrary to infinite order.

\begin{lemma}\label{lem:gap}

Let $\alpha$ be algebraic of infinite order and put $B=\lfloor\log(2n)/\lambda\rfloor+1$. For
each pair of distinct inputs there is a pair using only exponents in $[0,2nB]$ such that
exactly the same sign rows give equal answers.
\end{lemma}

\begin{proof}

Fix two inputs. Translate their entire finite set of occurring exponents so its minimum is
zero. This multiplies every row difference by the same nonzero scalar. Each row difference is
a signed sum of at most $2n$ powers. Split at consecutive used exponents $u<v$, and suppose
the low part is equal to the negative of the high part, with common nonzero magnitude at each
place. At an expanding place the low part is bounded by $|\alpha|_v^u$ times $2n$ at
archimedean places and times one elsewhere. At a contracting place use the high part and
$|\alpha|_v^v$ instead. At neutral places only the archimedean factor is needed.
Multiplication and the product formula give
\[
1\le 2n\,\exp[-(v-u)\lambda].
\]

Consequently a gap of at least $B=\lfloor\log(2n)/\lambda\rfloor+1$ splits a zero sum into two zero sums.

Replace every gap larger than $B$ by $B$. Inside each resulting band all exponents have
shifted together, so zero within that band is preserved. Splitting at all cuts, before and
after shortening, proves that exactly the same rows give equal answers. Distinct inputs stay
distinct because the exponent map is strictly increasing and two signed powers at distinct
exponents cannot coincide for a nontorsion radix. Every shortened exponent is at most $R=2nB$.
The shortening may depend on the pair, since its role is to certify collisions.
\end{proof}

\section{Cyclotomic quotients with large additive orders}\label{sec:projection}

Write $T=\mathbb Z[\alpha,\alpha^{-1}]$. Let $S$ contain the archimedean places and the finite
places with $|\alpha|_v\ne1$. Define $\|z\|_S=\prod_{v\in S}|z|_v^{w_v}$. Every nonzero $z\in
T$ satisfies $\|z\|_S\ge1$ by the product formula, since it has absolute value at most one
outside $S$. For positive integers $a$, $\|a\|_S\le a$. Also
$\|1\pm\alpha^r\|_S\le2e^{r\lambda}$ whenever the expression is nonzero and $r\ge0$.

\begin{lemma}\label{lem:proj}

Fix $d+1$ distinct odd primes. For every sufficiently large integer $h$, one can choose a
prime $p$ from this finite list such that, in the quotient
\[
T/\bigl(\Phi_p(\alpha^h)\bigr),\qquad \Phi_p(X)=1+X+\cdots+X^{p-1},
\]

the powers of $\alpha$ have period exactly $ph$, zero and both signs give $2ph+1$ different
letters, and the additive order of twice any nonzero letter difference is at least $c_\alpha
e^{h\lambda/2}$.
\end{lemma}

\begin{proof}

To choose $p$, consider the finite sets of primitive $p$-th roots on the unit circle. They are
disjoint for different primes. Choose disjoint neighborhoods of these sets. Each conjugate of
$\alpha$ on the unit circle, raised to the $h$-th power, rules out at most one prime. At most
$d$ primes are ruled out. For the chosen prime, $|\Phi_p(z)|$ is bounded below by a fixed
positive constant at every one of those conjugates raised to $h$, by compactness outside the
neighborhoods. At an expanding archimedean place the leading term has asymptotic ratio one to
the whole polynomial, and at a contracting archimedean place the value tends to one. At the
corresponding finite places the ultrametric inequality gives the exact absolute value from the
unique extreme term. The finite list of primes makes all bounds uniform in the choice. Thus
$D=\Phi_p(\alpha^h)$ satisfies
\[
\|D\|_S\ \ge\ c_\alpha e^{(p-1)h\lambda}.
\]

The identity $(\alpha^h-1)D=\alpha^{ph}-1$ gives a period dividing $ph$. The quotient is
additively generated by the first $ph$ powers, since the image of $\alpha$ is a unit of period
dividing $ph$. Tensoring with $\mathbb Q$ gives $E/(D)=0$. The quotient is therefore a
finitely generated torsion abelian group and is finite.

Reduce any two formal signed letters using their exponents modulo $ph$. After multiplying
their difference by a unit and possibly interchanging them, it is represented by $w=1$ or
$w=1\pm\alpha^r$ with $0\le r\le ph/2$, using the shorter circular difference. If the formal
letters differ this $w$ is nonzero in $E$. Suppose $2a w$ vanishes in the quotient for some
positive integer $a$. Then $2aw=Dz$ for a nonzero $z\in T$, and therefore
\[
c_\alpha e^{(p-1)h\lambda}
\ \le\ \|D\|_S\ \le\ \|2aw\|_S\ \le\ 4a e^{r\lambda}.
\]

Since $(p-1)h-r\ge h/2$, the claimed lower bound for $a$ follows after adjusting $c_\alpha$.
For large $h$ this excludes even equality of distinct formal letters, proving the exact period
and cardinality assertions.
\end{proof}

\section{Choosing one matrix for all projected differences}\label{sec:count}

\begin{proof}[Completion of the proof of Theorem~\ref{thm:class}]

Choose $H=\lceil C_\alpha\log n\rceil$ large enough that Lemma~\ref{lem:proj} gives additive
order greater than $n$ for every $h\in[H,2H]$ and that $2^H>R=2nB$. The least common multiple
of all the integers in this interval is at least $\binom{2H}{H}\ge2^H$. For the first
inequality, every positive integer up to $2H$ divides a member of the interval, and the
valuation of the central binomial coefficient at any prime $q$ is at most
$\lfloor\log_q(2H)\rfloor$, by the factorial valuation formula. The second inequality follows
by multiplying $(H+i)/i\ge2$ for $1\le i\le H$.

The projections jointly distinguish the shortened letters. For equal signs and differing
exponents in $[0,R]$, some $h$ fails to divide their difference, and hence $ph$ also fails to
divide it. Zero and opposite signs are distinguished in each projection. Each projection has
$O_\alpha(\log n)$ letters.

For completeness, draw all entries of $A$ as independent fair signs. Consider a nonzero vector
of projected differences with $k$ nonzero entries. The order bound ensures that changing the
sign on one nonzero entry changes a row answer, so the row collision probability is at most
$1/2$. Over all $h$, there are at most $\exp[O_\alpha(n/\log^2 n)]$ patterns with $k\le
n/\log^3 n$, by choosing the support and the ordered pair of letters in each position.
Hoeffding's inequality bounds the probability that one pattern vanishes on at least
$(1-2\eta)m$ rows by $\exp[-2(1/2-2\eta)^2m]$. Their total failure probability is summable in
$n$ for either sublinear rate in the table.

For larger supports, the number of available letter differences in a fixed projection is
$O_\alpha(\log^2 n)$. A particular nonzero difference $w$ occurs at least $\ell\gg_\alpha
n/\log^5 n$ times. Conditional on all other signs, these copies contribute $(2j-\ell)w$, where
$j$ is binomial with parameters $\ell,1/2$. All these values differ because the order of $2w$
exceeds $n$. By the central binomial bound, collision in one row has probability at most
$q_n=C_\alpha\log^{5/2}n/\sqrt n$. Over all projections there are at most $\exp(2n\log\log
n+O_\alpha(n))$ patterns. A fixed one agreeing on at least $(1-2\eta)m$ rows has probability
at most $2^m q_n^{(1-2\eta)m}$, by taking a union over row subsets. With any fixed
$\epsilon>0$ and
\[
m=\left\lceil\left(\frac4{1-2\eta}+\epsilon\right)\frac{n\log\log n}{\log n}\right\rceil ,
\]

the total dense failure probability is at most $\exp(-c_{\epsilon,\eta}n\log\log n)$ for
sufficiently large $n$. Combined with the sparse bound this is summable. No independence
between projections is needed.

Thus every nonzero projected pattern differs on more than $2\eta m$ rows. Every distinct
original pair has a differing projection after clipping. An unequal projected answer was
unequal before projection and before clipping. This gives disjoint correction balls and proves
the upper bound.

For a fixed finite alphabet in characteristic zero use the same counting argument directly in
its additive group. A nonzero difference has infinite additive order. The total number of
patterns is now $\exp(O(n))$, and the repeated difference occurs $\gg n/\log^3 n$ times for
larger supports. Thus $m=C_{\alpha,\eta} n/\log n$ suffices with summable failure probability.
This completes every row of the table.

Finally sample an infinite list of independent fair rows for each $n$, and use its prescribed
prefix in each case. Summability and the first Borel--Cantelli lemma give eventual success for
each fixed case with probability one. There are countably many algebraic radices up to their
rational minimal polynomial. Embedding the generated field elsewhere preserves all row
equalities. All characteristic-zero transcendental radices share the same ternary criterion,
and one prime-field estimate covers all extensions in each odd characteristic. The
intersection still has probability one, proving the universal array assertion.
\end{proof}

\section{Recovering a hidden binary wiring from inversion counts}\label{sec:wiring}

In the wiring problem, a box maps an $n$-bit word $z$ to the word whose bit $j$ is
$z_{p(j)}\oplus u_j$, for an unknown permutation $p$ and flip mask $u$. Bits are numbered from
least to most significant. A query chooses an additional output mask $a$ and returns the
number $I(a)$ of unordered input pairs whose unsigned order the masked box reverses.

\begin{proposition}\label{prop:wire}

Put $N=2^n$, let $r_j=\#\{i<j:p(i)<p(j)\}$ and set $c_j=(-1)^{u_j}2^{r_j}$. For each output mask $a$ the inversion count satisfies
\begin{equation}\label{eq:wire}
(N-1)-\frac{4I(a)}{N}=\sum_{j=0}^{n-1}(-1)^{a_j}c_j .
\end{equation}
The baseline count and any $n-1$ of the single-bit flip counts determine the wiring and all its flip bits.
\end{proposition}

\begin{proof}

For distinct words $x,y$, compare the signs of $y-x$ and of the corresponding difference
between the two masked outputs. The product of these signs is $+1$ for a concordant pair and
$-1$ for an inversion. Summing over ordered pairs gives $N(N-1)-4I(a)$, since an inversion
contributes twice. Now write $y=x\mathbin{\oplus}d$ and fix $d\ne0$. The input comparison uses
the highest set bit of $d$, and the output comparison uses the highest output whose input bit
is set in $d$. If the two comparisons use different wires, toggling the input bit of the first
wire in both words reverses exactly one comparison. The contribution summed over $x$ is zero.

If the two comparisons use the same wire $j$, then $d$ contains bit $p(j)$, and every other
set bit of $d$ belongs to a wire lower in both orders. There are $r_j$ such wires, hence
$2^{r_j}$ choices of $d$. Each gives the same contribution $N(-1)^{a_j+u_j}$. Summing over $j$
and dividing by $N$ proves \eqref{eq:wire}.

Write $I=I(0)$ and $I_j=I(e_j)$. Subtracting two instances of \eqref{eq:wire} gives
$c_j=2(I_j-I)/N$. The baseline gives $\sum_j c_j$, so one omitted coefficient is recovered by
subtraction. Its sign and magnitude, like those of every other coefficient, give $u_j$ and
$r_j$. Insert each successive output label $j$ at the zero-based position $r_j$ in a list.
After each insertion the list is ordered by the input positions of its labels. Thus the
position of $j$ in the final list is $p(j)$. Conversely every sequence with $0\le r_j\le j$
constructs a unique permutation this way.
\end{proof}

For a four-bit box, suppose the baseline count is 68. Flipping outputs 1, 2 and 3 gives 76, 36
and 84. Since $N=16$, the three differences give coefficients $1,-4,2$. The baseline says all
four coefficients add to $-2$, so the omitted coefficient is $-1$. The complete rank sequence
is $(0,0,2,1)$. The ordered outputs after the final insertion are $(1,3,0,2)$, giving
$p=(2,0,3,1)$. The negative signs mark outputs 0 and 2, so the hidden XOR mask is 5.

Let $F(n)=n\log\log n/\log n$. The leading constants below apply both to inversion counts and to the integer radix problem.

\begin{corollary}\label{cor:rate}

Fix $0\le\eta<1/4$. Let $M(n,\eta)$ denote the minimum number of fixed masks for the wiring
problem, or the minimum number of fixed sign rows for powers of an integer $b\ge2$ in
characteristic zero. In either case
\[
\frac{2}{1-2\eta}\ \le\
\liminf_{n\to\infty}\frac{M(n,\eta)}{F(n)}
\ \le\ \limsup_{n\to\infty}\frac{M(n,\eta)}{F(n)}
\ \le\ \frac{4}{1-2\eta}.
\]
Every deterministic adaptive noiseless protocol needs at least $(1-o(1))F(n)$ queries. The
supremum of the asymptotically correctable substitution fractions is $1/4$.
\end{corollary}

\begin{proof}

Every sign row corresponds to an output mask, and \eqref{eq:wire} is an invertible affine
change in each answer. Section~\ref{sec:count} therefore gives the upper constant for wirings
by taking $\alpha=2$, and for $b$ by taking $\alpha=b$.

Take $L=\lfloor\log_b n/(\log\log n)^2\rfloor$ for large $n$. Fix the first $L$ entries at
$1$. Independently allow each remaining exponent to take any of $0,\ldots,L-1$ and keep all
signs positive. At $b=2$ every one of these $L^{n-L}$ assignments is a wiring by the insertion
construction. Every signed sum is an integer of magnitude at most $n b^L=n^{1+o(1)}$.
Consequently each node of even an adaptive query tree has at most $2nb^L+1$ possible answers
on this family. Its leaves must distinguish $L^{n-L}$ inputs. Taking logs proves the
$(1-o(1))F(n)$ lower bound.

For fixed noiseless rows we can halve the leading logarithm of the number of outcomes per row
on a typical input. Give the free exponents the independent uniform distribution. For a row
sum $Y$, the bounded-sum inequality gives
\[
\Pr\!\left(\lvert Y-\mathbb EY\rvert>b^L\sqrt{2n\log n}\right)\le 2n^{-4}.
\]

For at most $n$ rows, at least half the inputs satisfy these bounds simultaneously for large
$n$. If there are more rows, the desired lower bound already holds. Each typical row has at
most $2b^L\sqrt{2n\log n}+2=n^{1/2+o(1)}$ possible integer values. Counting the transcripts
for at least $L^{n-L}/2$ inputs now gives $(2-o(1))F(n)$ as the minimum noiseless fixed
length. For an error-correcting schedule, deleting any $2\lfloor\eta m\rfloor$ rows leaves an
injective schedule. Otherwise two full transcripts would be within twice the correction radius
and could be changed into the same received word. This yields the lower constant
$2/(1-2\eta)$.

Use two coordinates of inputs in $\{-1,1\}^n$. These are all valid wiring coefficient vectors
at $b=2$ by taking all insertion ranks zero. There is a pair of inputs with difference
$(2,2,0,\ldots,0)$ and another pair with difference $(2,-2,0,\ldots,0)$. A sign row
distinguishes exactly one of these two pairs, according to whether its first two signs agree.
Their distances sum to $m$, so one is at most $\lfloor m/2\rfloor$. This proves the error
ceiling.
\end{proof}

\section*{AI Assistance Disclosure}

OpenAI's GPT-6.2 Astra was used in an autoresearch loop and reported this discovery to the author.

\begingroup\small
\begin{thebibliography}{9}
\bibitem{lindstrom}
B. Lindstr\"om, \emph{On a combinatorial problem in number theory}, Canadian Mathematical
Bulletin \textbf{8} (1965), 477--490.
\href{https://doi.org/10.4153/CMB-1965-034-2}{doi:10.4153/CMB-1965-034-2}.

\bibitem{emp}
G. Erdo\u{g}an, G. Maringer, and N. Polyanskii, \emph{Signature codes for a noisy adder
multiple access channel} (2022). \href{https://arxiv.org/abs/2206.10735}{arXiv:2206.10735}.

\bibitem{lenstra}
H. W. Lenstra, Jr., \emph{Finding small degree factors of lacunary polynomials}, in
\emph{Number Theory in Progress}, vol. 1, de Gruyter (1999), 267--276.
\href{https://doi.org/10.1515/9783110285581.267}{doi:10.1515/9783110285581.267}.

\bibitem{hoeffding}
W. Hoeffding, \emph{Probability inequalities for sums of bounded random variables}, Journal of
the American Statistical Association \textbf{58} (1963), 13--30.
\href{https://doi.org/10.1080/01621459.1963.10500830}{doi:10.1080/01621459.1963.10500830}.
\end{thebibliography}
\endgroup

\end{document}
